\documentclass[10pt,a4paper]{amsart}
\usepackage{amsmath,amssymb,amsthm,graphicx}
% Retained preamble for the Stacks Project chapter {CH} (shared project preamble).
% The chapter file pinned-original-source.tex only carries \input{{preamble}} on its first line,
% so no preamble text lives inside the chapter itself.  The Stacks Project uses one single shared
% preamble for every chapter: the snapshot stages it three times (sources/stacks-chow/preamble.tex,
% sources/stacks-etale-cohomology/preamble.tex and focused-seeds/AGP/sources/stacks-snapshot-20261011/
% preamble.tex) and those copies are identical in size.  The block below is copied mechanically,
% line for line, from sources/stacks-chow/preamble.tex lines 14-18 and 154-217: the Xy-pic setup and
% the native theorem declarations together with the mathematical macros.  Nothing was added, removed
% or redefined; the wrapper class and the amsmath/amssymb/amsthm/graphicx base packages are supplied
% by the assembler.
% For commutative diagrams we use Xy-pic
% Chapter: stacks-spaces-more-morphisms.tex -- More on Morphisms of Spaces.
\usepackage[all]{xy}

% We use 2cell for 2-commutative diagrams.
\xyoption{2cell}

% Theorem environments.
%
\theoremstyle{plain}
\newtheorem{theorem}[subsection]{Theorem}
\newtheorem{proposition}[subsection]{Proposition}
\newtheorem{lemma}[subsection]{Lemma}

\theoremstyle{definition}
\newtheorem{definition}[subsection]{Definition}
\newtheorem{example}[subsection]{Example}
\newtheorem{exercise}[subsection]{Exercise}
\newtheorem{situation}[subsection]{Situation}

\theoremstyle{remark}
\newtheorem{remark}[subsection]{Remark}
\newtheorem{remarks}[subsection]{Remarks}

\numberwithin{equation}{subsection}

% Macros
%
\def\lim{\mathop{\mathrm{lim}}\nolimits}
\def\colim{\mathop{\mathrm{colim}}\nolimits}
\def\Spec{\mathop{\mathrm{Spec}}}
\def\Hom{\mathop{\mathrm{Hom}}\nolimits}
\def\Ext{\mathop{\mathrm{Ext}}\nolimits}
\def\SheafHom{\mathop{\mathcal{H}\!\mathit{om}}\nolimits}
\def\SheafExt{\mathop{\mathcal{E}\!\mathit{xt}}\nolimits}
\def\Sch{\mathit{Sch}}
\def\Mor{\mathop{\mathrm{Mor}}\nolimits}
\def\Ob{\mathop{\mathrm{Ob}}\nolimits}
\def\Sh{\mathop{\mathit{Sh}}\nolimits}
\def\NL{\mathop{N\!L}\nolimits}
\def\CH{\mathop{\mathrm{CH}}\nolimits}
\def\proetale{{pro\text{-}\acute{e}tale}}
\def\etale{{\acute{e}tale}}
\def\QCoh{\mathit{QCoh}}
\def\Ker{\mathop{\mathrm{Ker}}}
\def\Im{\mathop{\mathrm{Im}}}
\def\Coker{\mathop{\mathrm{Coker}}}
\def\Coim{\mathop{\mathrm{Coim}}}

% Boxtimes
%
\DeclareMathSymbol{\boxtimes}{\mathbin}{AMSa}{"02}

%
% Macros for moduli stacks/spaces
%
\def\QCohstack{\mathcal{QC}\!\mathit{oh}}
\def\Cohstack{\mathcal{C}\!\mathit{oh}}
\def\Spacesstack{\mathcal{S}\!\mathit{paces}}
\def\Quotfunctor{\mathrm{Quot}}
\def\Hilbfunctor{\mathrm{Hilb}}
\def\Curvesstack{\mathcal{C}\!\mathit{urves}}
\def\Polarizedstack{\mathcal{P}\!\mathit{olarized}}
\def\Complexesstack{\mathcal{C}\!\mathit{omplexes}}
% \Pic is the operator that assigns to X its picard group, usage \Pic(X)
% \Picardstack_{X/B} denotes the Picard stack of X over B
% \Picardfunctor_{X/B} denotes the Picard functor of X over B
\def\Pic{\mathop{\mathrm{Pic}}\nolimits}
\def\Picardstack{\mathcal{P}\!\mathit{ic}}
\def\Picardfunctor{\mathrm{Pic}}
\def\Deformationcategory{\mathcal{D}\!\mathit{ef}}
\usepackage{cleveref}
\crefname{originallemma}{original lemma}{original lemmas}
\Crefname{originallemma}{Original lemma}{Original lemmas}
\makeatletter
\newlabel{lemma-algebraize-formal-closed-subscheme}{{08BG}{source}{}{sourcepointer-1}{}}
\newlabel{lemma-algebraize-formal-closed-subscheme@cref}{{[originallemma][0][]08BG}{[1][0][]source}{}{}{}}
\makeatother
\crefname{originallemma}{original lemma}{original lemmas}
\Crefname{originallemma}{Original lemma}{Original lemmas}
\makeatletter
\newlabel{lemma-quasi-finite-separated-pass-through-finite}{{082K}{source}{}{sourcepointer-2}{}}
\newlabel{lemma-quasi-finite-separated-pass-through-finite@cref}{{[originallemma][0][]082K}{[1][0][]source}{}{}{}}
\makeatother
\crefname{originalreference}{original reference}{original references}
\Crefname{originalreference}{Original reference}{Original references}
\makeatletter
\newlabel{spaces-morphisms-lemma-locally-finite-type-quasi-finite-part}{{04NW}{source}{}{sourcepointer-3}{}}
\newlabel{spaces-morphisms-lemma-locally-finite-type-quasi-finite-part@cref}{{[originalreference][0][]04NW}{[1][0][]source}{}{}{}}
\makeatother
\crefname{originalreference}{original reference}{original references}
\Crefname{originalreference}{Original reference}{Original references}
\makeatletter
\newlabel{spaces-morphisms-lemma-section-immersion}{{03KP}{source}{}{sourcepointer-4}{}}
\newlabel{spaces-morphisms-lemma-section-immersion@cref}{{[originalreference][0][]03KP}{[1][0][]source}{}{}{}}
\makeatother
\crefname{originalreference}{original reference}{original references}
\Crefname{originalreference}{Original reference}{Original references}
\makeatletter
\newlabel{spaces-morphisms-lemma-universally-closed-permanence}{{04NX}{source}{}{sourcepointer-5}{}}
\newlabel{spaces-morphisms-lemma-universally-closed-permanence@cref}{{[originalreference][0][]04NX}{[1][0][]source}{}{}{}}
\makeatother
\crefname{originaltheorem}{original theorem}{original theorems}
\Crefname{originaltheorem}{Original theorem}{Original theorems}
\makeatletter
\newlabel{theorem-grothendieck-existence}{{08BE}{source}{}{sourcepointer-6}{}}
\newlabel{theorem-grothendieck-existence@cref}{{[originaltheorem][0][]08BE}{[1][0][]source}{}{}{}}
\makeatother
\title{Algebraization of a compatible system of morphisms between formal algebraic spaces}
\author{The Stacks Project Authors}
\date{stacks-project:spaces-more-morphisms.tex:local-snapshot-20261011; GNU FDL 1.2 or later}
\begin{document}
\maketitle
\noindent\textit{Editorial note.} Counted span is the author's own proof environment, source lines 10261-10288, immediately following the statement at 10238-10259.\par\medskip
\noindent\textit{Editorial note.} The author ends with 'We omit the proof of uniqueness.'; that sentence is retained verbatim and the existence proof is complete inside the span.\par\medskip
\noindent\textit{Editorial note.} Three cross-chapter pointers (Morphisms of Spaces) plus three same-chapter statements are recorded in assembly.external\_references.\par\medskip
\noindent\textit{Allowed standard prerequisites (editorial).}\begin{itemize}
\item Algebraization of a formal closed subspace. chapter lemma lemma-algebraize-formal-closed-subscheme, cited in the proof. Produces the closed subspace Z inside X whose thickenings realise the given graphs.
\item Passing a quasi-finite morphism through a finite one. chapter lemma lemma-quasi-finite-separated-pass-through-finite, cited at proof lines 10280-10281. Combined with properness of Z over X, converts the quasi-finite p into a finite morphism (Zariski's main theorem).
\item Grothendieck's existence theorem. chapter theorem theorem-grothendieck-existence, cited at proof lines 10283-10284. Shows p\_*O\_Z = O\_X because this holds modulo I\textasciicircum{}n for all n, whence p is an isomorphism.
\item The quasi-finite part of a locally finite type morphism. Morphisms of Spaces, lemma spaces-morphisms-lemma-locally-finite-type-quasi-finite-part, cited at proof lines 10275-10276. Provides the open neighbourhood where p is quasi-finite, which properness then spreads to all of Z.
\end{itemize}\medskip
\ifdefined\hypertarget\hypertarget{sourcepointer-1}{}\fi
\noindent\textit{Source reference.} Original lemma 08BG: sources/stacks-spaces-more-morphisms/pinned-original-source.tex, source lines 10146--10146. Navigation pointer only: the counted proof cites this statement by \textbackslash{}ref as a standing input; the deduction the counted proof has to make is carried out inside the counted span, and no part of that deduction is delegated to this pointer.\par\medskip
\ifdefined\hypertarget\hypertarget{sourcepointer-2}{}\fi
\noindent\textit{Source reference.} Original lemma 082K: sources/stacks-spaces-more-morphisms/pinned-original-source.tex, source lines 7755--7755. Navigation pointer only: the counted proof cites this statement by \textbackslash{}ref as a standing input; the deduction the counted proof has to make is carried out inside the counted span, and no part of that deduction is delegated to this pointer.\par\medskip
\ifdefined\hypertarget\hypertarget{sourcepointer-3}{}\fi
\noindent\textit{Source reference.} Original reference 04NW: focused-seeds/geometry/S9-spaces-algebraize-morphism-spaces/ref-tags.txt, source lines 5--5. Navigation pointer only: the counted proof cites this statement by \textbackslash{}ref as a standing input; the deduction the counted proof has to make is carried out inside the counted span, and no part of that deduction is delegated to this pointer.\par\medskip
\ifdefined\hypertarget\hypertarget{sourcepointer-4}{}\fi
\noindent\textit{Source reference.} Original reference 03KP: focused-seeds/geometry/S9-spaces-algebraize-morphism-spaces/ref-tags.txt, source lines 6--6. Navigation pointer only: the counted proof cites this statement by \textbackslash{}ref as a standing input; the deduction the counted proof has to make is carried out inside the counted span, and no part of that deduction is delegated to this pointer.\par\medskip
\ifdefined\hypertarget\hypertarget{sourcepointer-5}{}\fi
\noindent\textit{Source reference.} Original reference 04NX: focused-seeds/geometry/S9-spaces-algebraize-morphism-spaces/ref-tags.txt, source lines 7--7. Navigation pointer only: the counted proof cites this statement by \textbackslash{}ref as a standing input; the deduction the counted proof has to make is carried out inside the counted span, and no part of that deduction is delegated to this pointer.\par\medskip
\ifdefined\hypertarget\hypertarget{sourcepointer-6}{}\fi
\noindent\textit{Source reference.} Original theorem 08BE: sources/stacks-spaces-more-morphisms/pinned-original-source.tex, source lines 9990--9990. Navigation pointer only: the counted proof cites this statement by \textbackslash{}ref as a standing input; the deduction the counted proof has to make is carried out inside the counted span, and no part of that deduction is delegated to this pointer.\par\medskip
\begin{lemma}
\label{lemma-algebraize-morphism}
Let $A$ be a Noetherian ring complete with respect to an ideal $I$.
Write $S = \Spec(A)$ and $S_n = \Spec(A/I^n)$. Let $X$, $Y$ be algebraic
spaces over $S$. For $n \geq 1$ we set $X_n = X \times_S S_n$ and
$Y_n = Y \times_S S_n$. Suppose given a compatible system of
commutative diagrams
$$
\xymatrix{
& & X_{n + 1} \ar[rd] \ar[rr]_{g_{n + 1}} & & Y_{n + 1} \ar[ld] \\
X_n \ar[rru] \ar[rd] \ar[rr]_{g_n} & & Y_n \ar[rru] \ar[ld] & S_{n + 1} \\
& S_n \ar[rru]
}
$$
Assume that
\begin{enumerate}
\item $X \to S$ is proper, and
\item $Y \to S$ is separated of finite type.
\end{enumerate}
Then there exists a unique morphism of algebraic spaces $g : X \to Y$
over $S$ such that $g_n$ is the base change of $g$ to $S_n$.
\end{lemma}

\begin{proof}
The morphisms $(1, g_n) : X_n \to X_n \times_S Y_n$ are closed immersions
because $Y_n \to S_n$ is separated
(Morphisms of Spaces, Lemma \ref{spaces-morphisms-lemma-section-immersion}).
Thus by Lemma \ref{lemma-algebraize-formal-closed-subscheme}
there exists a closed subspace $Z \subset X \times_S Y$
proper over $S$ whose base change to $S_n$ recovers
$X_n \subset X_n \times_S Y_n$. The first projection $p : Z \to X$
is a proper morphism (as $Z$ is proper over $S$, see
Morphisms of Spaces, Lemma
\ref{spaces-morphisms-lemma-universally-closed-permanence})
whose base change to $S_n$ is an isomorphism for all $n$.
In particular, $p : Z \to X$ is quasi-finite on an open subspace
of $Z$ containing every point of $Z_0$ for example by
Morphisms of Spaces, Lemma
\ref{spaces-morphisms-lemma-locally-finite-type-quasi-finite-part}.
As $Z$ is proper over $S$ this open neighbourhood is all of $Z$.
We conclude that $p : Z \to X$ is finite by Zariski's main theorem
(for example apply
Lemma \ref{lemma-quasi-finite-separated-pass-through-finite}
and use properness of $Z$ over $X$ to see that the immersion is
a closed immersion). Applying the equivalence of
Theorem \ref{theorem-grothendieck-existence}
we see that $p_*\mathcal{O}_Z = \mathcal{O}_X$ as this is true
modulo $I^n$ for all $n$. Hence $p$ is an isomorphism and we obtain
the morphism $g$ as the composition $X \cong Z \to Y$.
We omit the proof of uniqueness.
\end{proof}

\end{document}
